\begin{table}[htbp]
\centering
\caption{Efeito do orçamento de parada. \emph{Gap}: distância percentual média ao melhor valor conhecido, nas 28 instâncias de treino. \emph{Tempo}: mediana por execução, medida sem concorrência em seis instâncias-amostra (uma por família e tipo). O valor adotado ($K=800$) está em negrito. O gap decresce ao longo de toda a faixa medida, sem platô --- razão pela qual $K$ é declarado como orçamento computacional, e não derivado de um ponto de convergência.}\label{tab:orcamento}
\begin{tabular}{r rrr rrr}
\toprule
& \multicolumn{3}{c}{\emph{gap} médio (\%)} & \multicolumn{3}{c}{tempo mediano (s)} \\
\cmidrule(lr){2-4}\cmidrule(lr){5-7}
$K$ & GRASP & Reativo & Tabu & GRASP & Reativo & Tabu \\
\midrule
50 & 3{,}49 & 3{,}11 & 5{,}99 & --- & --- & --- \\
100 & 3{,}15 & 2{,}71 & 5{,}47 & --- & --- & --- \\
200 & 2{,}39 & 2{,}27 & 4{,}68 & 4{,}6 & 6{,}9 & 1{,}6 \\
400 & 2{,}17 & 1{,}93 & 4{,}28 & 9{,}5 & 11{,}8 & 3{,}9 \\
800 & \textbf{1{,}86} & \textbf{1{,}65} & \textbf{4{,}04} & \textbf{15{,}6} & \textbf{22{,}7} & \textbf{6{,}8} \\
1600 & 1{,}67 & 1{,}52 & 3{,}64 & 27{,}9 & 52{,}5 & 13{,}9 \\
3200 & 1{,}47 & 1{,}20 & 3{,}28 & --- & --- & --- \\
\bottomrule
\end{tabular}
\end{table}
